\begin{abstract}

% Recent advancements in GPU-accelerated, photo-realistic simulators have enabled a new data-scaling path for robot learning. By using massive visual randomization, these simulators enable generalization to open-world scenes. Using a teacher-student-bootstrap approach, we target the challenging humanoid loco-manipulation task of opening diverse doors, which remain unsolved in the open-world, RGB-only settings. We propose an exploration scheme based on staged-reset to stabilize teacher policy training, and a novel GRPO-based fine-tuning strategy tailored for Sim2Real RL. The resulting policy, \method{}, trained only on synthetic simulation data, can zero-shot open various doors (with lever handles, push-bars, fixed handles, etc.) in open-world scenes, outperforming human teleoperators by 31.7\% in task completion time. To the best of our knowledge, it is the first end-to-end door-opening policy ever deployed on a humanoid robot.

Recent progress in GPU-accelerated, photorealistic simulation has opened a scalable data-generation path for robot learning, where massive physics and visual randomization allow policies to generalize beyond curated environments. Building on these advances, we develop a teacher-student-bootstrap learning framework for vision-based humanoid loco-manipulation, using articulated-object interaction as a representative high-difficulty benchmark. Our approach introduces a staged-reset exploration strategy that stabilizes long-horizon privileged-policy training, and a GRPO-based fine-tuning procedure that mitigates partial observability and improves closed-loop consistency in sim-to-real RL. Trained entirely on simulation data, the resulting policy achieves robust zero-shot performance across diverse door types and outperforms human teleoperators by up to 31.7\% in task completion time under the same whole-body control stack. This represents the first humanoid sim-to-real policy capable of diverse articulated loco-manipulation using pure RGB perception.

\end{abstract}